\section{有限位数楕円境界の中心商に必要な正則性の証明}
\label{CSBGsec:pair-finite-order-background}
この節は固定した周囲の正則な中心商を得るための証明を収める。
一般中心商の分類や全中心の同定は別リポジトリCayley--Sklyanin論文の対象である。
以下で$B$は標数零の体上の左右noetherian connected graded AS正則代数、
$W\in Z(B)$は正次数の正則斉次元、$A=B/(W)$とする。
$\mathcal P_A=\operatorname{thick}\{\pi A(n):n\in\mathbb Z\}$、
$\mathcal Q_A=\mathrm D^b(\qgr A)/\mathcal P_A$と置く。
有限graded自由行列因子化と次数付き特異圏の比較には
\cite[Theorems~4.7, 5.5]{CSBGCCKM}を用いる。
右加群には反対環を適用し、中心元のnormalizing automorphismは恒等である。
従って有限syzygy後のcoherent対象を二周期因子化で扱える。
以下の補題の引用入力と仮定を保ち、一般の非正則底へ拡張しない。
\begin{lemma}[有限射影因子化の縮約と分母払い]
\label{CSBGlem:gluing-contraction}
$2$が可逆な左右noetherian環$S$と中心正則元$w$を取り、
有限生成射影加群の行列因子化を$d^2=w$と書く。
$S/(w)$上のcokernelが有限射影次元なら、この因子化は縮約可能である。
さらに中心正則元$u$で局所化した後に縮約可能なら、
ある$N$について$u^N\mathrm{id}=dK+Kd$となる。
次数付きの場合も、因子化の次数shiftを含めて同じ主張が成立する。
\end{lemma}
\begin{proof}
$\bar S=S/(w)$上の二周期分解から、cokernelの有限射影次元は
各周期syzygyの射影性を意味する。従って分解は分裂し、
$\bar S$上で$d\bar H+\bar Hd=1$となる奇写像を取れる。
有限射影性からこれを$S$上の$H$へ持ち上げると
$dH+Hd=1+wT$となる。
左辺は$d$と可換なので、$w$の正則性により$[d,T]=0$。
ゆえに$H'=H-\frac12dT$は$dH'+H'd=1$を満たす。
局所化後の$H'$の有限個の成分の分母を$u^N$で払い、
$S\hookrightarrow S[u^{-1}]$の単射性を使えば最後の等式を得る。
\end{proof}

\begin{lemma}[有限大域次元からtwistsによる生成へ]
\label{CSBGlem:global-finite-to-perfect}
$A$は次数1生成で、$\operatorname{qgr}A$が有限大域次元を持つなら、
$\mathrm D^b(\operatorname{qgr}A)=\mathcal P_A$である。
\end{lemma}
\begin{proof}
大域次元を$e$とし、graded自由表示を繰り返して
$0\to K\to L_e\to\cdots\to L_0\to M\to0$を取る。
各$L_i$は$\pi A(n)$の有限直和である。
最後のYoneda拡大は$\operatorname{Ext}^{e+1}(M,K)=0$なので、
$M$は有限複体$[L_e\to\cdots\to L_0]$の直和因子である。
line object自体が射影的であるとは仮定していない。
\end{proof}

\begin{proposition}[中心零化元による十分条件]
\label{CSBGprop:global-annihilators}
$u,v\in Z(B)_1$が$A$上正則列をなし、
$\mathrm D^b(\operatorname{qgr}A/(u,v))$が構造対象のtwistsで厚く生成されるとする。
正次数の斉次中心元$p_1,\ldots,p_r\in Z(A)$が
次数付き特異圏の全対象を零化し、ある$N$について
$(u,v)^N\subset(p_1,\ldots,p_r)A$なら、$\mathcal Q_A=0$である。
\end{proposition}
\begin{proof}
graded lift $M$のKoszul複体$K(p_1,\ldots,p_r;M)$のcohomologyは
$(p_1,\ldots,p_r)$に殺され、従って$(u,v)$の冪に殺される。
有限filtrationの各商は$A/(u,v)$加群である。
正則列のKoszul分解から$\pi(A/(u,v))(n)\in\mathcal P_A$なので、
境界の生成性によりこのKoszul複体は$\mathcal P_A$に入る。
一方、$\mathrm D^b(\operatorname{gr}A)\to\mathcal Q_A$は
次数付きperfect複体を殺すから次数付き特異圏を経由する。
従って各$p_i$は$\mathcal Q_A$で零作用し、反復coneは
$M$のshiftとtwistの直和へ分裂する。ゆえに$M$自身も零となる。
\end{proof}

\begin{proposition}[二つの中心chartの貼合せ]
\label{CSBGprop:global-two-charts}
$u,v\in Z(B)_1$が$B$上正則で$A$上正則列をなし、
$A[u^{-1}]_0,A[v^{-1}]_0$が有限大域次元を持つとする。
$\operatorname{qgr}A/(u,v)$の有界導来圏が構造対象のtwistsで厚く生成されるなら、
$\mathcal Q_A=0$である。
\end{proposition}
\begin{proof}
各coherent対象を行列因子化で表し、各chart上で縮約する。
補題\ref{CSBGlem:gluing-contraction}の分母払いから、対象ごとに選んだ
$u^r,v^s$が$\mathcal Q_A$で零作用する。
二重Koszul coneのcohomologyは$(u^r,v^s)$に殺され、
$(u,v)^{r+s-1}\subset(u^r,v^s)$から境界の有限filtrationを持つ。
前命題と同じ議論でconeはperfectである。
中心作用は自然変換なので、最初の零coneの直和分解とも可換し、
二重coneは元の対象の四つのshift・twistの直和となる。
従って元の対象は零である。一様な$r,s$は要求しない。
\end{proof}

\begin{lemma}[perfect性とSerre双対性からの正確な次元]
\label{CSBGlem:gluing-serre-dimension}
$A$は左右noetherian connected graded AS--Gorensteinで、
balanced dualizing complexが$A_\nu(-\ell)[d+1]$（$d\ge0$）で、
$\pi A\ne0$とする。
$\mathrm D^b(\operatorname{qgr}A)=\mathcal P_A$なら
$\operatorname{gldim}\operatorname{qgr}A=d$である。
\end{lemma}
\begin{proof}
\cite[Theorem~4.2, Corollary~4.3]{CSBGYZSerre}は左右の$\chi$と
有限cohomological dimensionを与える。
特に$\pi A(n)$は有限射影次元を持つ。
Chen\cite[Theorem~3.10]{CSBGChenGeneralizedSerre}の一般化Serre双対性は
有限射影次元の対象に適用でき、その関手は
$S(M)=M_\nu(-\ell)[d]$となる。
仮定により全coherent対象が適用範囲に入るので、
\[
 \operatorname{Ext}^i(M,N)^*
 \simeq\operatorname{Ext}^{d-i}(N,M_\nu(-\ell)).
\]
$i>d$では右辺が負次数ゆえ零。
また$M=\pi A$、$N=M_\nu(-\ell)$と取ると
$\operatorname{Ext}^d(M,N)^*\simeq\operatorname{Hom}(N,N)\ne0$なので等号となる。
\end{proof}

\begin{lemma}[超越的中心値の有限大域次元]
\label{CSBGlem:global-transcendental}
$k_0\subset\CC$、$R_0$は$k_0$-代数、$z\in Z(R_0)$は$k_0$上超越的とする。
$R=R_0\otimes_{k_0}\CC$が左右noetherianで$\operatorname{gldim}R\le n$なら、
$k_0$上超越的な全ての$\lambda\in\CC$について
$\operatorname{gldim}R/(z-\lambda)\le n$である。
\end{lemma}
\begin{proof}
$\lambda$を含む$\CC/k_0$の超越基底$T$を選び、$K=k_0(T\setminus\{\lambda\})$とする。
体拡大は自由かつ忠実平坦で、$K$-線形な$\CC\to K$の$1\mapsto1$という分裂がある。
従って大域次元の上界とnoetherian性は$R_K=R_0\otimes K$へ降下する。
$z$の冪は$K$上でも独立なので、中心局所化
$R_\eta=R_K\otimes_{K[z]}K(z)$も大域次元高々$n$。
$z\mapsto\lambda$で$K(z)=K(\lambda)$とみなすと
$R_\eta\otimes_{K(z)}\CC=R/(z-\lambda)$である。
最後の体拡大は代数的分離拡大で、右辺は$R$の商ゆえnoetherian。
\cite[Lemmas~3.2, 3.3]{CSBGRSVBertini}の仮定が揃い、上界が保たれる。
任意の超越的体拡大が有限大域次元を保つと仮定したわけではない。
\end{proof}

\begin{theorem}[固定した良いambientでのvery-generalな正則性]
\label{CSBGthm:global-very-general}
$B$は中心一次元$u,v$を持つ左右noetherian AS正則五次元代数とし、
$G\in Z(B)_3$、$\operatorname{gr}_{x,y,z}B=S\otimes\CC[u,v]$とする。
$S$は整域、$\operatorname{gr}G=\gamma\in Z(S)_3$は非零正則元とする。
全ての$f\in\CC[u,v]_3$について$B/(G+f)$は$\CC[u,v]$上自由であり、
$\operatorname{qgr}(S/(\gamma))$は滑らかな楕円曲線のcoherent圏と仮定する。
有限表示の定義体$k_0$に対し、$a,d$が$k_0(b,c)$上代数的独立なら
\[
 \operatorname{gldim}\operatorname{qgr}
 B/(G+au^3+bu^2v+cuv^2+dv^3)=3.
\]
\end{theorem}
\begin{proof}
$D_u=B[u^{-1}]_0$はnoetherianで有限大域次元を持つ。
実際、中心局所化は大域次元を増やさず、
$B[u^{-1}]=D_u[u,u^{-1}]$から制限・誘導と評価$u=1$で上界を降下できる。
$t=v/u$とすると商の$u$-chartは
$D_u/(G/u^3+bt+ct^2+dt^3+a)$。
括弧内の$a$を除く中心元は$k_0(b,c,d)$上定義され、
最高主変数成分が非零正則なので、その冪は独立である。
前補題でchartの有限大域次元を得る。
$v$-chartも$d$を中心値として同様に扱える。
命題\ref{CSBGprop:global-two-charts}と
補題\ref{CSBGlem:gluing-serre-dimension}が結論を与える。
\end{proof}

\begin{lemma}[被約基底上の中心の持上げ]
\label{CSBGlem:global-center-lifting}
$H=B(E,\mathcal L,\tau)$は標数零の滑らかな楕円曲線のTHCRで、
$\mathcal L$はample、$\tau$は有限位数のtranslationとする。
$R$を被約有限型可換代数とし、filtered $R$-代数$U$が
$\operatorname{gr}U=H\otimes R$、各$F_dU$は有限自由で、
閉点への基底変更と両立するとする。
このとき各$z\in Z(H)_d$は$Z(U)\cap F_dU$へ持ち上がり、
$U$は有限型可換中心部分代数上有限となる。
\end{lemma}
\begin{proof}
斉次生成元の持上げ$u_i$を取り、有限自由加群間の交換子写像
$M_d:F_dU\to\bigoplus_iF_{d+\deg u_i}U$を作る。
Rains\cite[Theorem~2.13]{CSBGRainsElliptic}により各閉点$s$で
$\operatorname{gr}Z(U_s)=Z(H)$なので、$M_d$のrankは一定である。
被約性により高すぎるminorは恒等的に零となり、
最大rankのminorで被覆すれば$\ker M_d$は有限射影で、形成は基底変更と可換する。
最高次symbolの写像$\ker M_d\to Z(H)_d\otimes R$は
各閉点で全射、従ってNakayamaにより全射である。
右辺が自由なので、各斉次中心元は大域的に持ち上がる。
$H$は有限位数translationのnorm line bundleによる中心上有限である。
中心の有限生成元と中心上の有限加群生成元を持ち上げ、次数を下げる帰納法で
$U$も持ち上げ中心部分代数上有限と分かる。
ここではRainsの体上の定理を非被約基底へ拡張していない。
\end{proof}

\begin{theorem}[有限位数境界を持つ固定ambientの開な正則範囲]
\label{CSBGthm:global-open}
定理\ref{CSBGthm:global-very-general}の良い固定ambientを取り、さらに
$H=S/(\gamma)=B(E,\mathcal L,\tau)$のtranslationが有限位数とする。
ある非零多項式$P(a,b,c,d)$が存在して、$P\ne0$上では
中心三次商のtails大域次元が3となる。
周囲のPBW/Koszul解消とnoetherian AS条件が基底変更と両立する構成では、
この十分条件は任意の代数閉拡大体上でも成立する。
特に周囲の表示が有理数係数なら、有理数四係数の正則な商がZariski稠密に存在する。
\end{theorem}
\begin{proof}
$D_u=B[u^{-1}]_0$、$t=v/u$、$G_u=G/u^3$とする。
全係数を動かす$u$-chartの全代数は
\[
 \mathcal U=D_u[a,b,c,d]/(G_u+a+bt+ct^2+dt^3)
 \simeq D_u[b,c,d]
\]
なので左右noetherianかつ有限大域次元である。
$R=\CC[t,a,b,c,d]$をfiltration次数零とすると
$\operatorname{gr}\mathcal U=H\otimes R$で、各filtered pieceは有限自由。
前補題から$\mathcal U$は有限型可換中心部分代数$C_u$上有限である。
この$C_u$は係数環$R_0=\CC[a,b,c,d]$上も有限型である。
ここで\cite[Proposition~2.1, proof steps (6)--(7)]{CSBGRSVBertini}を
$(R,S,A)=(R_0,C_u,\mathcal U)$へ適用する。
その証明はある非零$P_u\in R_0$を逆にした上で、$\mathcal U$が$R_0$上平坦で、
相対対角が有限射影解消を持つことを与える。
従って全ての幾何学的fiberで有限大域次元が成立し、体拡大も許される。
$v$-chartでも$d$を消去して$P_v$を得る。
$P=P_uP_v$とし、二chartの貼合せとSerre双対性で結論する。
$\mathbb Q^4$は複素四次元affine空間でもZariski稠密なので、
非空開集合には有理数係数点が稠密に存在する。
\end{proof}

\begin{lemma}[位数6の具体的な楕円曲線計算]\label{CSBGlem:order-six}
$E_1=V(x^3+y^3+z^3+xyz)$で$O=[1:-1:0]$を原点とすると、
$p=[1:-1:-1]$は位数6であり、$2p=[1:0:-1]$である。
\end{lemma}
\begin{proof}
$E_1$は滑らかで、$O$は変曲点である。
$p$の接線$2x+y+z=0$へ代入すると
$7(y+z)(y-z)^2/8$となり、第三交点は$[0:1:-1]$である。
この点と$O,[1:0:-1]$は共線なので、接線の第三交点を反転すると表示した$2p$を得る。
$2p$は$O$と異なる変曲点なので位数3、$p$自体はHessianが非零で変曲点ではない。
従って$p$の位数は6である。各代入は現行の厳密検算に含まれる。
\end{proof}
$h=-1$の場合は$x,y$の交換が$S_{-1}$と$S_1$の次数付き同型を与えるので、境界translationの位数も同じである。
